% FILE: CSE-19_CT-1.tex
\documentclass[11pt, a4paper]{article}
\usepackage[a4paper, top=2.5cm, bottom=2.5cm, left=2cm, right=2cm]{geometry}
\usepackage{fontspec}
\usepackage[english]{babel}
\usepackage{amsmath}
\usepackage{amssymb}
\usepackage{enumitem}

\babelfont{rm}{Noto Sans}

\begin{document}

%%%%%%%%%%%%%%%%%%%%%%%%%%%%%%%%%%%%%%%%%%%%%%%%%%  
% QUESTION_START  
%%%%%%%%%%%%%%%%%%%%%%%%%%%%%%%%%%%%%%%%%%%%%%%%%%
% id=cse19-ct1-seca-q1  
% discipline=CSE  
% year=19  
% exam_type=CT  
% section=A  
% ct_number=1  
% question_number=1  
% topic=Applications of Derivatives  
% subtopic=Maxima \& Minima  
% difficulty=Medium  
% length=Medium  
% question_type=New  
% frequency=1  
% appearances=  
% OCR_STATUS=VERIFIED

\section*{CT-1 (Sec A) - Q1}

Question:
Locate \& describe the Relative extrema of $f(x) = 3x^{5/3} - 15x^{2/3}$.

Solution:

\begin{description}
    \item[Step 1: Find the first derivative $f'(x)$] \ \\
    The given function is:
    \[ f(x) = 3x^{5/3} - 15x^{2/3} \]
    Differentiating with respect to $x$:
    \[ f'(x) = 3 \cdot \frac{5}{3}x^{2/3} - 15 \cdot \frac{2}{3}x^{-1/3} \]
    \[ f'(x) = 5x^{2/3} - 10x^{-1/3} \]
    Factoring out $5x^{-1/3}$:
    \[ f'(x) = 5x^{-1/3}(x - 2) = \frac{5(x - 2)}{x^{1/3}} \]

    \item[Step 2: Identify the critical points] \ \\
    Critical points occur where $f'(x) = 0$ or where $f'(x)$ is undefined:
    \begin{itemize}
        \item $f'(x) = 0 \implies 5(x - 2) = 0 \implies x = 2$.
        \item $f'(x)$ is undefined at $x = 0$ (since the denominator becomes zero, though $f(0)$ is defined).
    \end{itemize}
    Thus, the critical points are $x = 0$ and $x = 2$.

    \item[Step 3: Analyze intervals of increase and decrease (First Derivative Test)] \ \\
    We test the sign of $f'(x) = \frac{5(x - 2)}{x^{1/3}}$ in the intervals divided by $x = 0$ and $x = 2$:
    \begin{enumerate}
        \item \textbf{Interval $(-\infty, 0)$:} \\
        Choose a test point $x = -1$:
        \[ f'(-1) = \frac{5(-1 - 2)}{(-1)^{1/3}} = \frac{-15}{-1} = 15 > 0 \implies f(x) \text{ is increasing.} \]

        \item \textbf{Interval $(0, 2)$:} \\
        Choose a test point $x = 1$:
        \[ f'(1) = \frac{5(1 - 2)}{(1)^{1/3}} = \frac{-5}{1} = -5 < 0 \implies f(x) \text{ is decreasing.} \]

        \item \textbf{Interval $(2, \infty)$:} \\
        Choose a test point $x = 3$:
        \[ f'(3) = \frac{5(3 - 2)}{(3)^{1/3}} = \frac{5}{3^{1/3}} > 0 \implies f(x) \text{ is increasing.} \]
    \end{enumerate}

    \item[Step 4: Describe the Relative Extrema] \ 
    \begin{itemize}
        \item \textbf{At $x = 0$:} The derivative changes from positive to negative. Therefore, $f(x)$ has a \textbf{relative maximum} at $x = 0$. \\
        The relative maximum value is:
        \[ f(0) = 3(0)^{5/3} - 15(0)^{2/3} = 0 \]
        
        \item \textbf{At $x = 2$:} The derivative changes from negative to positive. Therefore, $f(x)$ has a \textbf{relative minimum} at $x = 2$. \\
        The relative minimum value is:
        \[ f(2) = 3(2)^{5/3} - 15(2)^{2/3} = 3 \cdot 2 \cdot 2^{2/3} - 15 \cdot 2^{2/3} = -9 \cdot 2^{2/3} \]
    \end{itemize}
\end{description}

Final Answer:
\[ 
\boxed{\text{Relative Maximum at } (0, 0), \quad \text{Relative Minimum at } (2, -9\cdot 2^{2/3})} 
\]

%%%%%%%%%%%%%%%%%%%%%%%%%%%%%%%%%%%%%%%%%%%%%%%%%%  
% QUESTION_END  
%%%%%%%%%%%%%%%%%%%%%%%%%%%%%%%%%%%%%%%%%%%%%%%%%%

\end{document}